% ===== Document Preamble =====
\documentclass[11pt]{article}
\usepackage[margin=0.5in]{geometry}
\usepackage{amsmath,amssymb}
\usepackage{xcolor}
\usepackage{enumitem}
\usepackage{needspace}
\usepackage{comment}

% ===== Colors =====
\definecolor{AnsRed}{RGB}{170,0,0}
\definecolor{WrongBlue}{RGB}{0,0,160}
\definecolor{SkipGray}{RGB}{120,120,120}

% ===== Choice list styling (tight, consistent) =====
\setlist[enumerate]{%
  leftmargin=1.5em,
  itemsep=0.25em,
  topsep=0.25em,
  label=\textbf{\Alph*.}
}

% ===== Style Wrappers =====
\newcommand{\QuestionStyle}[1]{{\large\bfseries\boldmath #1}}
\newcommand{\AnswerStyle}[1]{{\large\bfseries\boldmath #1}}
\newcommand{\ReasonStyle}[1]{{\normalsize #1}}

% ===== Show / Hide Reasoning Toggle =====
\newif\ifshowreason
\showreasontrue      % Instructor version
% \showreasonfalse   % Student version

% ===== Choice Item Macros (Answer + Reasoning) =====
% Usage:
%   \Correct{Answer}{Reasoning}
%   \Wrong{Answer}{Reasoning}

\newcommand{\Correct}[2]{%
  \item \AnswerStyle{\textcolor{AnsRed}{#1}}%
  \ifshowreason\\[-0.2em]
    \ReasonStyle{\textcolor{AnsRed}{#2}}%
  \fi
}

\newcommand{\Wrong}[2]{%
  \item \AnswerStyle{\textcolor{WrongBlue}{#1}}%
  \ifshowreason\\[-0.2em]
    \ReasonStyle{\textcolor{WrongBlue}{#2}}%
  \fi
}

% ===== Multiple-Choice Question Block Macro =====
% Usage: \QBlock{Q1}{stem...}{<items>}{answer key}

\newcommand{\QBlock}[4]{%
  \Needspace{8\baselineskip}%
  \noindent
  \begin{minipage}{\linewidth}
    \QuestionStyle{#1.~#2}

    \vspace{0.25em}
    \begin{enumerate}
      #3
    \end{enumerate}

    \vspace{0.1em}
    \noindent\textcolor{AnsRed}{\textit{\textbf{Answer:}~#4}}

    \vspace{0.25em}
    \hrule height 0.4pt
  \end{minipage}
  \vspace{0.75em}
}

% ===== Open-Ended Question Block Macro =====
% Usage: \OEBlock{Q1}{stem...}{answer text}

\newcommand{\OEBlock}[3]{%
  \Needspace{5\baselineskip}%
  \noindent
  \begin{minipage}{\linewidth}
    \QuestionStyle{#1.~#2}

    \vspace{0.25em}
    \noindent\textcolor{AnsRed}{\textit{\textbf{Answer:}~#3}}

    \vspace{0.25em}
    \hrule height 0.4pt
  \end{minipage}
  \vspace{0.75em}
}

% ===== Multi-Part Question Block Macro =====
% Usage: \MPBlock{Q1}{stem...}{<parts>}

\newcommand{\MPBlock}[3]{%
  \Needspace{8\baselineskip}%
  \noindent
  \begin{minipage}{\linewidth}
    \QuestionStyle{#1.~#2}

    \vspace{0.25em}
    #3

    \vspace{0.25em}
    \hrule height 0.4pt
  \end{minipage}
  \vspace{0.75em}
}

% ===== MC Sub-Part Macro =====
% Usage: \MCPart{a}{stem...}{<items>}{answer key}

\newcommand{\MCPart}[4]{%
  \par\vspace{0.35em}
  \noindent\textbf{(#1)}~#2
  \begin{enumerate}
    #3
  \end{enumerate}
  \vspace{0.1em}
  \noindent\textit{\textbf{Answer:}~#4}
  \par
}

% ===== Open-Ended Sub-Part Macro =====
% Usage: \OEPart{a}{stem...}{answer text}

\newcommand{\OEPart}[3]{%
  \par\vspace{0.35em}
  \noindent\textbf{(#1)}~#2
  \par\vspace{0.15em}
  \noindent\textit{\textbf{Answer:}~#3}
  \par
}

% ===== Skipped Question Note Macro =====
% Usage: \SkipBlock{Q1}{reason}

\newcommand{\SkipBlock}[2]{%
  \noindent\textcolor{SkipGray}{\textit{[#1 omitted: #2]}}%
  \vspace{0.5em}\par
}

% ===== Document Start =====

\begin{document}

\begin{center}
{\huge \textbf{Physics: Electricity and Magnetism}}\\[0.25em]
{\Large Isomorphic Practice Questions}\\[0.25em]
{\large Questions generated below. Multiple-choice answers in RED; incorrect answers in BLUE.}
\end{center}

\vspace{0.75em}

\QBlock{Q1}{Two point charges are in vacuum. A positive charge $+q$ is at $x=-0.30\,\text{m}$ on the $x$-axis, and a negative charge $-q$ is at the origin. Using Coulomb's law in the form $\mathbf F_{\text{on }2}=k\,q_1q_2\,\hat{\mathbf r}_{1\to2}/r^2$ (where $\hat{\mathbf r}_{1\to2}$ points from charge 1 to charge 2), what is the direction of the Coulomb force on the negative charge, and what is the direction of $\hat{\mathbf r}_{+q\to -q}$?}{
  \Wrong{Force on $-q$ points in the $+x$ direction; $\hat{\mathbf r}_{+q\to -q}$ points in the $+x$ direction}{This treats the interaction as repulsive; opposite charges attract, so the force on $-q$ must point toward $+q$ (to the left).}
  \Wrong{Force on $-q$ points in the $+x$ direction; $\hat{\mathbf r}_{+q\to -q}$ points in the $-x$ direction}{This gets both directions wrong: $\hat{\mathbf r}_{+q\to -q}$ should point from $x=-0.30\,\text{m}$ to $0$, i.e. $+x$, and the force on $-q$ is attractive toward $-x$.}
  \Wrong{Force on $-q$ points in the $-x$ direction; $\hat{\mathbf r}_{+q\to -q}$ points in the $-x$ direction}{The force direction is correct, but the unit vector is defined from $+q$ (left) to $-q$ (at the origin), which is in the $+x$ direction, not $-x$.}
  \Correct{Force on $-q$ points in the $-x$ direction; $\hat{\mathbf r}_{+q\to -q}$ points in the $+x$ direction}{Because the charges are opposite, the force on $-q$ is attractive and points toward $+q$, i.e. in the $-x$ direction. The separation unit vector $\hat{\mathbf r}_{+q\to -q}$ points from $x=-0.30\,\text{m}$ to $0$, which is $+\hat{\mathbf i}$; since $q_1q_2<0$ in $\mathbf F=kq_1q_2\hat{\mathbf r}/r^2$, the force is opposite $\hat{\mathbf r}$, giving $-x$.}
}{D}

\QBlock{Q2}{A positive point charge $+Q$ is fixed at the origin. A negative point charge $-Q$ is placed at the point $(x,y)=(-0.20\,\text{m},0)$ on the negative $x$-axis. Which statement correctly gives (i) the direction of the electric field at that location due to $+Q$ and (ii) the direction of the electric force on the negative charge placed there?}{
  \Wrong{Electric field in the $+x$ direction; electric force in the $-x$ direction}{This reverses the electric field direction: the field of $+Q$ points away from the origin, which at $x<0$ is toward $-x$.}
  \Wrong{Electric field in the $-x$ direction; electric force in the $-x$ direction}{It correctly sets $\mathbf E$ toward $-x$, but forgets that $\mathbf F=q\mathbf E$ reverses direction for a negative charge ($q<0$).}
  \Correct{Electric field in the $-x$ direction; electric force in the $+x$ direction}{The electric field of a positive source charge points radially outward from the origin; at a point on the negative $x$-axis, that is the $-x$ direction. The force is $\mathbf F=q\mathbf E$; with $q=-Q$, the force points opposite $\mathbf E$, so it is in the $+x$ direction.}
  \Wrong{Electric field in the $+y$ direction; electric force in the $-y$ direction}{This confuses the field direction with an unrelated axis; the field from a point charge at the origin is radial along the line connecting the origin to the point (here the $x$-axis).}
}{C}

\QBlock{Q3}{A point charge is located at the geometric center of a regular tetrahedron that is used as a Gaussian surface (the charge is exactly equidistant from all four faces by symmetry). Which statement is correct?}{
  \Wrong{Electric flux is zero through each face}{A nonzero enclosed charge implies total flux $\Phi_{\text{tot}}=Q_{\text{enc}}/\varepsilon_0\neq 0$, so the flux cannot be zero through every face.}
  \Correct{Electric flux is the same through each face}{By symmetry, all four faces are equivalent when the charge is at the tetrahedron's center, so each face must carry the same flux. Gauss's law gives $\Phi_{\text{tot}}=Q_{\text{enc}}/\varepsilon_0$, hence $\Phi_{\text{face}}=\Phi_{\text{tot}}/4=Q_{\text{enc}}/(4\varepsilon_0)$.}
  \Wrong{The magnitude of the electric field is the same everywhere on each face}{Equal flux does not require uniform $|\mathbf E|$; points on a face are at varying distances from the charge, so $|\mathbf E|$ varies across the face.}
  \Wrong{None of the above is correct}{Option B follows directly from symmetry plus Gauss's law, so ``none of the above'' is incorrect.}
}{B}

\QBlock{Q4}{An infinitely long straight wire carries a uniform (constant) linear charge density. A closed cylindrical Gaussian surface is chosen coaxial with the wire. Which part(s) of the Gaussian surface have nonzero electric flux?}{
  \Wrong{Both the curved side surface and the two flat end caps}{This incorrectly assumes the field has a component through the end caps; for an infinite line charge $\mathbf E$ is purely radial, parallel to the end caps, so flux there is zero.}
  \Correct{Only the curved side surface}{For an infinite line charge, the electric field is radial, so it is perpendicular to the curved side's area vector and gives $\mathbf E\cdot d\mathbf A\neq 0$ on the side. On the end caps, $d\mathbf A$ is along the cylinder axis while $\mathbf E$ is radial (tangent to the caps), so $\mathbf E\cdot d\mathbf A=0$ there.}
  \Wrong{Only the two flat circular end caps}{This would require an axial electric field component, but an infinite line charge produces no axial field component by symmetry.}
  \Wrong{No part of the surface}{A closed surface enclosing nonzero charge must have net flux $Q_{\text{enc}}/\varepsilon_0\neq 0$, so some part of the surface must contribute (the curved side does).}
}{B}

\QBlock{Q5}{A positive point charge $+3.0\,\mu\text{C}$ and a negative point charge $-3.0\,\mu\text{C}$ are separated by $0.60\,\text{m}$ in vacuum. The charges are then moved closer until they are only $0.20\,\text{m}$ apart. How does the electric potential energy $U$ of the two-charge system change?}{
  \Wrong{It increases (becomes more positive or less negative)}{For opposite charges $U=kq_1q_2/r$ is negative, and decreasing $r$ makes $U$ more negative, not less negative.}
  \Correct{It decreases (becomes less positive or more negative)}{The potential energy of two point charges is $U=\dfrac{kq_1q_2}{r}$. Here $q_1q_2<0$, so $U$ is negative; reducing $r$ (from $0.60$ m to $0.20$ m) increases $|U|$ and makes $U$ more negative, so the energy decreases.}
  \Wrong{It becomes zero}{$U=0$ only in the limit $r\to\infty$ (for the conventional choice $U(\infty)=0$), not at a finite separation.}
  \Wrong{It remains unchanged (but not zero)}{Because $U\propto 1/r$, changing the separation necessarily changes the potential energy; it is not distance-independent.}
}{B}

\QBlock{Q6}{A small test charge is moved from point $A$ to point $B$ in a region of uniform electric field $\mathbf E$. Which factor determines the electric potential difference $\Delta V=V_B-V_A$ between the two points?}{
  \Wrong{The path taken between $A$ and $B$}{Electrostatic fields are conservative, so $\Delta V$ is path-independent; only endpoints matter.}
  \Wrong{The time it takes to move from $A$ to $B$}{Potential difference is defined from the field and geometry, not from the motion timing of the particle.}
  \Wrong{The mass of the test charge}{$\Delta V$ is a property of the field and the two points; it does not depend on the particle's mass.}
  \Correct{The component of the displacement from $A$ to $B$ along $\mathbf E$}{In a uniform field, $\Delta V=-\int_A^B \mathbf E\cdot d\mathbf l=-\mathbf E\cdot\Delta\mathbf r$. Therefore only the projection of $\Delta\mathbf r$ along $\mathbf E$ affects $\Delta V$; any displacement perpendicular to $\mathbf E$ gives zero contribution to $\mathbf E\cdot d\mathbf l$.}
}{D}

\QBlock{Q7}{A parallel-plate capacitor (vacuum between plates) has plate area $A$ and separation $d$. If the plate separation is increased to $3d$ while the plate area remains $A$, the capacitance is:}{
  \Wrong{Tripled}{This incorrectly assumes $C$ grows with separation; for parallel plates $C\propto 1/d$.}
  \Wrong{Unchanged}{$C$ depends on geometry; changing $d$ changes $C$ through $C=\varepsilon_0A/d$.}
  \Correct{Reduced by a factor of 3}{For a parallel-plate capacitor, $C=\dfrac{\varepsilon_0A}{d}$. If $d\to 3d$, then $C\to \dfrac{\varepsilon_0A}{3d}=\dfrac{C}{3}$, so it is reduced by a factor of 3.}
  \Wrong{Reduced by a factor of 9}{A factor of 9 would require $d$ to increase by a factor of 9 (not 3); $C$ scales linearly with $1/d$, not $1/d^2$.}
}{C}

\QBlock{Q8}{Two point charges are separated by $16.0\,\text{cm}$ and have charges $+1.50\,\mu\text{C}$ and $-1.50\,\mu\text{C}$, respectively. What is the magnitude of the electric field at the point midway between them? (Use $k=8.99\times10^9\,\text{N}\cdot\text{m}^2/\text{C}^2$.)}{
  \Correct{$4.21\times10^6\,\text{N/C}$}{At the midpoint, the distance to each charge is $r=0.080\,\text{m}$, so each produces magnitude $E_1=\dfrac{k|q|}{r^2}$. The fields point the same way (away from $+$ and toward $-$), so they add: $E=2\dfrac{k|q|}{r^2}=2\dfrac{(8.99\times10^9)(1.50\times10^{-6})}{(0.080)^2}\approx4.21\times10^6\,\text{N/C}$.}
  \Wrong{$2.11\times10^6\,\text{N/C}$}{This is the field from only one charge at the midpoint; it misses that the second charge's field points in the same direction and must be added.}
  \Wrong{$8.43\times10^6\,\text{N/C}$}{This overcounts the field (e.g., effectively doubling twice or using an incorrect smaller distance than $0.080\,\text{m}$ in the $1/r^2$ law).}
  \Wrong{$0$}{Opposite charges do not cancel the field at the midpoint; cancellation at the midpoint occurs for equal like charges (fields oppose), not for a dipole (fields align between charges).}
}{A}

\QBlock{Q9}{A total charge of $+4.20\,\mu\text{C}$ is uniformly distributed throughout an insulating sphere of radius $0.100\,\text{m}$. What is the magnitude and direction of the electric field at a point located $0.250\,\text{m}$ from the surface of the sphere?}{
  \Wrong{$5.39\times10^5\,\text{N/C}$, radially outward}{This magnitude is too large and typically comes from using the wrong distance from the center (for example, using $0.250\,\text{m}$ instead of $0.350\,\text{m}$ in $E=kQ/r^2$).}
  \Correct{$3.08\times10^5\,\text{N/C}$, radially outward}{The point is outside the sphere, so the field equals that of a point charge at the center by Gauss's law. The distance from the center is $r=0.100+0.250=0.350\,\text{m}$, so $E=\dfrac{k|Q|}{r^2}=\dfrac{(8.99\times10^9)(4.20\times10^{-6})}{(0.350)^2}\approx3.08\times10^5\,\text{N/C}$, directed radially outward because $Q>0$.}
  \Wrong{$3.08\times10^5\,\text{N/C}$, radially inward}{The magnitude matches using the correct $r$, but the direction is wrong: for a positive total charge, the electric field points outward.}
  \Wrong{$1.54\times10^5\,\text{N/C}$, radially outward}{This is about half the correct value and can result from an algebra error with the inverse-square dependence (e.g., treating $E\propto 1/r$ instead of $1/r^2$).}
}{B}

\QBlock{Q10}{A particle with charge $-Q$ is placed at the center of an initially uncharged conducting hollow sphere. After electrostatic equilibrium is reached, what charges are induced on the inner and outer surfaces of the conductor?}{
  \Wrong{Inner: $-Q$; outer: $+Q$}{This has the signs reversed: the inner surface must carry charge opposite the central charge to ensure $\mathbf E=0$ in the conducting material.}
  \Wrong{Inner: $+Q$; outer: $0$}{If $+Q$ appears on the inner surface, charge conservation for an initially neutral conductor requires $-Q$ to appear on the outer surface, not $0$.}
  \Correct{Inner: $+Q$; outer: $-Q$}{The electric field inside the conducting material must be zero at equilibrium, so the inner surface must induce charge $+Q$ to balance the enclosed $-Q$ (so that a Gaussian surface in the metal encloses zero net charge). Because the conductor started neutral, the remaining induced charge must be $-Q$ on the outer surface so the conductor's total charge remains zero.}
  \Wrong{Inner: $0$; outer: $+Q$}{The inner surface cannot remain uncharged when a charge is inside the cavity; otherwise there would be an electric field within the conductor, violating electrostatic equilibrium.}
}{C}

\QBlock{Q11}{The electric potential along the $x$-axis is given by $V(x)=4x^3-6x$, where $V$ is in volts and $x$ is in meters. What is the electric field $\mathbf E(x)$?}{
  \Wrong{$\mathbf E=(12x^2-6)\,\hat{\mathbf i}$}{This is $dV/dx$, but the electric field is $\mathbf E=-\nabla V$, so the sign should be reversed.}
  \Correct{$\mathbf E=(-12x^2+6)\,\hat{\mathbf i}$}{In one dimension, $E_x=-\dfrac{dV}{dx}$. Since $\dfrac{dV}{dx}=12x^2-6$, we get $E_x=-(12x^2-6)=-12x^2+6$, so $\mathbf E=(-12x^2+6)\hat{\mathbf i}$.}
  \Wrong{$\mathbf E=(-4x^4+6x^2)\,\hat{\mathbf i}$}{This comes from incorrectly integrating or using the wrong power rule; differentiating $4x^3$ should give $12x^2$, not an $x^4$ term.}
  \Wrong{$\mathbf E=(4x^3-6x)\,\hat{\mathbf i}$}{This incorrectly assumes $\mathbf E=V\hat{\mathbf i}$ rather than using $\mathbf E=-\nabla V$.}
}{B}

\QBlock{Q12}{A $12\,\text{V}$ battery is connected in series with a $5.0\,\mu\text{F}$ capacitor, a $3.0\,\Omega$ resistor, an ammeter, and a switch that is initially open. A voltmeter is connected in parallel across the capacitor. At the instant the switch is closed ($t=0^+$), what are the ammeter and voltmeter readings, respectively?}{
  \Correct{$4.0\,\text{A},\ 0\,\text{V}$}{At $t=0^+$ the capacitor is initially uncharged, so $V_C(0^+)=0$ and it behaves like a short circuit. Therefore the full battery voltage is across the resistor and $I(0^+)=V/R=12\,\text{V}/3.0\,\Omega=4.0\,\text{A}$, while the voltmeter across the capacitor reads $0\,\text{V}$.}
  \Wrong{$0\,\text{A},\ 12\,\text{V}$}{This corresponds to the long-time steady state after charging, not the instant just after closing when the capacitor voltage is still zero.}
  \Wrong{$0\,\text{A},\ 0\,\text{V}$}{While $V_C(0^+)=0$ is correct, the current is not zero because an uncharged capacitor initially allows current (acts like a short).}
  \Wrong{$4.0\,\text{A},\ 12\,\text{V}$}{If the capacitor had $12\,\text{V}$ across it at that instant, the resistor would have $0\,\text{V}$ and the current would be $0\,\text{A}$; this mixes up initial and final behavior.}
}{A}

\QBlock{Q13}{A $9.0\,\text{V}$ battery (negligible internal resistance) charges a $470\,\mu\text{F}$ capacitor through a $1.5\,\text{k}\Omega$ resistor. How long does it take for the capacitor's charge to reach $70\%$ of its maximum value?}{
  \Wrong{$0.25\,\text{s}$}{This is too small compared to the time constant $RC=(1.5\times10^3)(470\times10^{-6})\approx0.705\,\text{s}$; reaching $70\%$ requires a time on the order of $RC$.}
  \Wrong{$0.71\,\text{s}$}{This incorrectly assumes $70\%$ occurs at exactly one time constant; in fact $q(RC)/Q_{\max}=1-e^{-1}\approx63\%$, not $70\%$.}
  \Wrong{$1.20\,\text{s}$}{This is too large and typically comes from using the wrong logarithm (for example using $-\ln(0.70)$ instead of $-\ln(0.30)$).}
  \Correct{$0.85\,\text{s}$}{For charging, $q(t)=Q_{\max}(1-e^{-t/RC})$. Setting $q/Q_{\max}=0.70$ gives $e^{-t/RC}=0.30$ so $t=-RC\ln(0.30)$; with $RC=(1.5\times10^3)(470\times10^{-6})\approx0.705\,\text{s}$, $t\approx0.705\times1.204\approx0.85\,\text{s}$.}
}{D}

\QBlock{Q14}{At one instant an electron ($q=-1.60\times10^{-19}\,\text{C}$) is moving in the $xy$-plane with velocity components $v_x=-4.0\times10^5\,\text{m/s}$ and $v_y=+6.0\times10^5\,\text{m/s}$. A magnetic field of magnitude $0.50\,\text{T}$ points in the positive $y$ direction. At that instant, the magnetic force on the electron is:}{
  \Wrong{$(3.2\times10^{-14}\,\text{N})\,(-\hat{\mathbf k})$}{This has the right magnitude but the wrong direction; it effectively treats the charge as positive or mishandles the sign of $v_x$.}
  \Wrong{$(6.4\times10^{-14}\,\text{N})\,(+\hat{\mathbf k})$}{This overestimates the magnitude (e.g., by using $v=\sqrt{v_x^2+v_y^2}$ instead of the perpendicular component $|v_x|$ to $\mathbf B$).}
  \Correct{$(3.2\times10^{-14}\,\text{N})\,(+\hat{\mathbf k})$}{The magnetic force is $\mathbf F=q\,\mathbf v\times\mathbf B$. With $\mathbf B=B\hat{\mathbf j}$ and $\mathbf v=v_x\hat{\mathbf i}+v_y\hat{\mathbf j}$, we have $\mathbf v\times\mathbf B=v_xB\,\hat{\mathbf k}$; since $v_x<0$, $\mathbf v\times\mathbf B$ points in $-\hat{\mathbf k}$, and multiplying by the electron's negative charge reverses it to $+\hat{\mathbf k}$. The magnitude is $|F|=|q||v_x|B=(1.60\times10^{-19})(4.0\times10^5)(0.50)=3.2\times10^{-14}\,\text{N}$.}
  \Wrong{$(3.8\times10^{-14}\,\text{N})\,(+\hat{\mathbf k})$}{This is a arithmetic/magnitude error; the correct magnitude from $|q||v_x|B$ is $3.2\times10^{-14}\,\text{N}$.}
}{C}

\end{document}
